% Generated draft: requires mathematical and source review.
\documentclass[11pt]{article}
\usepackage[margin=25mm]{geometry}
\usepackage{fontspec}
\setmainfont{Latin Modern Roman}
\usepackage{amsmath,amssymb,mathtools,unicode-math}
\setmathfont{Latin Modern Math}
\usepackage{enumitem,needspace,xurl,hyperref}
\setlength{\parindent}{0pt}
\setlength{\parskip}{5pt}
\setlength{\emergencystretch}{3em}
\newcommand{\N}{\mathbb{N}}
\newcommand{\Z}{\mathbb{Z}}
\newcommand{\Q}{\mathbb{Q}}
\newcommand{\R}{\mathbb{R}}
\newcommand{\C}{\mathbb{C}}
\newcommand{\F}{\mathbb{F}}
\newcommand{\A}{\mathbb{A}}
\newcommand{\PP}{\mathbb{P}}
\newcommand{\E}{\mathbb{E}}
\newcommand{\eps}{\varepsilon}
\newcommand{\dd}{\,\mathrm{d}}
\begin{document}
\section*{UnsolvedMath ID: 2307012}
\textbf{Source problem code:} catalogue entry
\subsection{Definitions and mathematical statement}
The source asks for higher local integrability of derivatives of an \(n\)-dimensional \(K\)-``quasi-analytic'' function, without defining that historical term. If it means a quasiregular map in the usual bounded-distortion sense, a possible precise setting is \(n\ge2\), \(K\ge1\), a domain \(\Omega\subseteq\R^n\), and a continuous map \(f:\Omega\to\R^n\) with \(f\in W^{1,n}_{\mathrm{loc}}\) and
\[
\|Df(x)\|^n\le KJ_f(x)\quad\text{for almost every }x,\qquad
J_f=\det Df,
\]
where \(\|Df\|\) is operator norm. This is explicitly a conditional modern interpretation, not an identification supplied by the excerpt.

With that meaning and interpreting the printed ``\(L\)-integrable'' as \(L^p\)-integrable, the requested assertion is
\[
\forall n\ge2\ \forall K\ge1\ \exists c=c(K,n)>0\quad
\forall f\text{ in the intended class}\quad
\forall p\in[n,n+c],\quad
\partial_j f_i\in L^p_{\mathrm{loc}}(\Omega)\quad(1\le i,j\le n).
\]
Local integrability means \(\int_U|\partial_j f_i|^p\,\dd x<\infty\) for every relatively compact \(U\subset\Omega\). The gain \(c\) must be independent of the particular map and domain. No injectivity hypothesis is imposed in the question.

The source credits Gehring with the injective case. The supplied OPEN-TRIAGE summary describes the general noninjective case as open in the 2018 edition and says no definitive resolution was located. Neither the definition nor a precise current theorem citation is supplied, so this status claim is retained as metadata rather than an independently verified statement about the modern quasiregular class. The distortion convention, dimension range, function-versus-map interpretation, and missing \(p\) exponent all require source review.
\subsection{Short English statement}
Prove a uniform positive gain beyond exponent \(n\) in the local integrability of derivatives of the intended bounded-distortion maps, without injectivity. The historical class and the printed integrability notation are undefined.
\subsection{Sources}
\begin{itemize}
\item Local UnsolvedMath registry, permanent ID \(2307012\), exact source code \texttt{catalogue entry}, ``Research Problems in Function Theory --- Problem 7.12.''
\item Hayman, \emph{Research Problems in Function Theory} (2018), Problem 7.12, attributed to F. W. Gehring and crediting him with the injective case; \url{https://arxiv.org/abs/1809.07200}. Source-TeX extraction; recorded access 2026-07-29.
\item Supplied OPEN-TRIAGE summary concerns general noninjective quasi-analytic functions but supplies no complete reference for either the known theorem or the claimed remaining status.
\end{itemize}
\end{document}
